\documentclass[conference]{IEEEtran}
\IEEEoverridecommandlockouts
\usepackage{amsmath,amssymb,amsfonts}
\usepackage{booktabs}
\usepackage{tabularx}
\usepackage{graphicx}
\usepackage{xcolor}
\usepackage{cite}
\usepackage{microtype}

\title{Mathematical Formulation of the Proactive Golden-Hour Bridge Model (PGBM) for Post-Earthquake UAV-Assisted Response}

\author{\IEEEauthorblockN{Tamim Dewan || Faiaz Mahmud}
\IEEEauthorblockA{\textit{Department of Computer Science and Engineering}\\
\textit{University Name}\\
Email: author@email.com}
}

\begin{document}

\maketitle

\begin{abstract}
This document presents the mathematical formulation of the Proactive Golden-Hour Bridge Model (PGBM).
\end{abstract}

\section{Introduction}
\subsection{Problem in One Sentence}
After Scout UAVs detect survivors, the Command Center must decide at each event time $t$ which available Responder UAVs should execute which feasible tours---ordered sequences $\pi_j=(i_1,\dots,i_k)$---to maximize expected lifesaving value under battery, payload-item, and 3D reachability constraints, before ground rescue arrives.

\section{System Entities}

\subsection{Survivors}
Tasks $i$ with $p_i\in\mathbb{R}^3$, $t_i^{\text{det}}$, $c_i\in[0,1]$, $\sigma_i\in[0,1]$, $r_i^q\in\mathbb{Z}_{\ge0}$ for $q\in\{\text{med},\text{water},\text{food}\}$.

\subsection{Responder UAVs}
Platforms $j$ with $p_j$, $p_j^{\text{base}}\in\mathbb{R}^3$, $B_j(t)\in[0,B_{\max}]$, capacity $K$, inventory $I_j^q$, payload $m_j=\sum_q I_j^q w_q$, status $\in\{\text{Idle},\text{EnRoute},\text{Returning}\}$.

\subsection{Command Center}
% To be filled later

\subsection{Environment}
$\mathcal{E}=(V,O)$ with $V\subset\mathbb{R}^3$ free space, $O\subset\mathbb{R}^3$ static obstacles, $V\cap O=\emptyset$, hazard $R(p)$ deferred to $R=0$ in V1.

\section{System Model}
% Define entities, environment, tasks

\section{Problem Formulation}
% Offloading + Routing

\section{Conclusion}
% Summary

\end{document}
